
6.3.
Pour toute @@M0@@-algèbre graduée @@M1@@, on désigne par @@M2@@ la @@M3@@-algèbre @@M4@@, graduée de façon évidente. D'après 1.13, on définit un homomorphisme fonctoriel : @@M5@@ en donnant sa valeur pour @@M6@@, soit
@@M7@@

Cet homomorphisme est compatible avec la multiplication définie en 6.1 sur @@M8@@~: il suffit d'après 1.13 de la montrer pour les éléments @@M9@@ de @@M10@@~; or on a :
@@M11@@

Notons @@M12@@ l'endomorphisme (fonctoriel par rapport à @@M13@@) de @@M14@@, défini par :
@@M15@@
Alors, @@M16@@,
@@M17@@
En effet, @@M18@@ et @@M19@@ sont deux homomorphismes fonctoriels de groupes abéliens, et il suffit pour qu'ils soient égaux qu'ils coïncident sur @@M20@@, ce qui est évident.

On en déduit que @@M21@@ commute aux opérations de @@M22@@~: en effet, il suffit là encore de le montrer pour @@M23@@, et on a :
@@M24@@
